\chapter{作品概述}
\section{问题背景与应用场景}
企业智能体需要读取工单、告警、邮件、代码和知识库，再调用文件、网络或代码托管工具完成任务。外部资料和工具元数据与可信任务要求共同进入模型上下文，攻击者因而可以通过资料中的指令诱导智能体越权。OWASP将提示注入列为LLM应用风险之一\cite{owasp}；MCP规范也将工具调用的用户控制、数据访问和工具安全纳入宿主责任\cite{mcp}。

本作品选择安全研判作为主要场景。例如，用户要求分析一份含有危险命令的告警，系统需要完整引用该命令并说明影响，同时限制其触发凭据读取或网络外发。这里需要同时回答两个问题：报告是否正确保留证据，以及工具是否产生未经授权的实际效果。只对文本作恶意分类，或仅检查模型最终回答，都不足以回答这两个问题。

溯界将授权判断放在实际工具执行之前。模型可以提出动作，可信宿主依据任务、来源和当前状态决定动作能否执行；检测结果用于补充风险解释。现有网络和终端防护可继续作为部署防线，溯界补充其任务授权和上下文来源维度。

\section{相关工作与作品切入点}
已有研究形成了几条互补路线。CaMeL通过能力和信息流约束区分控制流与不可信数据\cite{camel}；Task Shield核验指令和工具调用是否服务用户目标\cite{taskshield}；AgentSpec以领域策略描述并执行运行约束\cite{agentspec}。AuthGraph比较授权与实际来源关系\cite{authgraph}，AgentVisor研究语义权限分离与违规后的反馈\cite{agentvisor}，ContainmentBench和AID-Guard分别关注暴露后的遏制评估和委派效果的状态化授权\cite{containment,aidguard}。MCP-SafetyBench与MCPTox提供MCP安全和工具投毒的基准研究\cite{mcpsafety,mcptox}。

溯界的切入点是将任务授权、来源用途和实际效果收据放入同一执行链，支持危险资料研判、跨角色委派和跨会话记忆。表\ref{tab:position}说明本作品使用已有思路时增加的具体工程对象和验收方式，机制维度的横向对照见表\ref{tab:compare}。具体实现和验证范围分别见第二、三章。

\begin{table}[htbp]\centering
\caption{作品切入点与验收对象}\label{tab:position}
\begin{tabularx}{\textwidth}{P{0.24\textwidth}YY}\toprule
问题 & 工程对象 & 验收对象\\\midrule
资料影响工具授权 & 任务契约、Grant、来源用途信封 & 相同资料下的允许/拒绝动作与真实效果\\
委派和记忆传播 & 子权限、来源父链、版本原文与CAS & 实际主体/会话、原历史、版本与回执\\
审批后动作漂移 & 参数和版本绑定、执行前复核 & 指定commit/ref或载荷、确认与实际提交\\
安全与业务判定 & 独立业务与副作用Oracle & 实际报告事实、读取记录、Receiver与Git状态\\\bottomrule
\end{tabularx}
\end{table}

\section{功能定位与实用性}
原型提供13类业务配置：告警报告、公开IOC传输、网页核查、构件校验、Git检视、Git发布、多层委派、记忆复核、记忆更新、事件时间线、Issue分级、附件一致性核对和知识版本协调。每类配置对应一种业务入口；配置数量不等于独立测试任务数。

面向大型企事业单位，系统可作为智能体平台接入真实工具前的授权与审计组件：把一次操作关联到任务目标、输入来源、批准对象和实际结果。当前验证使用合成资料、封闭接收端和本地Git远端，便于在不触及真实业务资产的情况下演示和复核。

表\ref{tab:rules}按赛题列明的规则逐条给出实现状态与证据位置。表中的“实例已验证”表示已有受控工具链用例，“部分实现”表示尚有明确缺项；两者均不代表对任意载荷的完整覆盖。

\begin{table}[htbp]\centering
\caption{赛题列明规则与实现状态}\label{tab:rules}
\begin{tabularx}{\textwidth}{P{0.215\textwidth}P{0.135\textwidth}Y}\toprule
赛题列明规则 & 状态 & 实现对象与证据位置\\\midrule
\texttt{fs.write} 写凭据、Shell启动文件、定时任务、VCS钩子与IDE设置 & 实例已验证 & 不可覆盖敏感名单，放宽Grant下仍拒绝（\ref{sec:sensitive}节，真实stdio MCP测试）\\
\texttt{net.fetch} 外发会话内已读secret & 实例已验证 & 来源用途信封与目的地判定，封闭Receiver给出实际接收记录（\ref{sec:effect}节）\\
\texttt{exec} 执行\texttt{base64}解码或管道下载脚本 & 实例已验证 & 仅允许已注册程序及确切参数，拒绝未登记命令；未提供通用Shell安全执行（\ref{sec:effect}节）\\
代码检视期间\texttt{git.push} & 实例已验证 & 检视保持远端ref不变；发布绑定确切commit/remote/ref并需审批（\ref{sec:approval}节）\\
子Agent违规传递secret & 实例已验证 & 子Grant收窄资源/操作/预算，父Host接收保留来源的结果（\ref{sec:delegation}节）\\
放行/审批/阻断三态 & 实例已验证 & 策略DSL输出ALLOW/ASK/BLOCK，硬拒绝优先（\ref{sec:effect}节）\\
运行时风险等级 & 部分实现 & 决策记录含由实际决策反推的风险字段；尚无独立校准的风险分级器\\\bottomrule
\end{tabularx}
\end{table}

\section{威胁模型与工作假设}
攻击者可控制用户输入中的越权文本、网页和文档正文、MCP工具描述或返回、记忆正文，以及被劫持子Agent的输出。模型输出与这些资料均作为待核验数据。可信边界包括认证与审批入口、任务契约与Grant、网关、来源存储、已注册效果适配器、受控执行器和宿主环境。

系统保护的目标包括：阻止秘密流向未授权目的地，保持任务和权限边界，在传播与变换后保留来源用途约束，并记录关键动作和实际效果。可用性以合法业务结果另行判断，不能由阻断次数替代。宿主失陷、任意远程工具内部行为、跨崩溃一致性以及通用操作系统隔离不在当前验收范围；具体实施范围见第二、三章。
